\noindent
<<school>> \hfill NOM : \underline{\hspace{5.5cm}} \\
<<course_title>> <<course>><<group_paren>> \hfill <<date_long>> \\
Professeur : <<professor>> \hfill Durée : <<duration>>

\begin{center}
    \huge\underline{\textbf{Examen \#<<exam_number>> --- Version <<version>>}}<<solution_banner>>
\end{center}

\vspace{0.3cm}
\hrule
\vspace{0.3cm}

\textbf{CONSIGNES POUR L'EXAMEN :} les éléments suivants sont exigés dans les solutions de problèmes :
\begin{itemize}[leftmargin=*, noitemsep]
    \item schéma de la situation avec identification des variables ;
    \item données initiales ;
    \item exposé clair de chaque étape de la solution avec titre ou courte phrase explicative, équations utilisées, résolution algébrique (voir ci-dessous) et calculs numériques (avec unités).
    \begin{itemize}
        \item Isolez algébriquement la variable recherchée avant de remplacer par les valeurs numériques.
        \item Faites le moins d'approximations possible lors de vos calculs.
        \item Donnez les réponses finales avec au moins trois chiffres significatifs.
    \end{itemize}
\end{itemize}
Les calculatrices programmables, téléphones cellulaires et autres appareils multimédias sont interdits.\\
Lorsque vous aurez terminé, brochez le questionnaire sur le dessus de votre examen.

\vspace{0.3cm}
\hrule
\vspace{0.5cm}

\textbf{Cet examen comporte <<summary>> pour un total de <<total>> points.
\\
\\
Note importante : Si vous ne réussissez pas à trouver la réponse à une partie d’une
question et vous avez besoin de cette réponse pour les parties suivantes, vous pouvez
inventer une valeur pour poursuivre le problème et vous ne serez pas pénalisé.e tant
que votre démarche reste cohérente à la valeur inventée!}
\vspace{5cm}
\begin{center}
    \Huge \textbf{\rule{1cm}{0.4pt}/<<total>> points}
\end{center}
\newpage
